\textbf{1.} D'après le théorème de Pythagore, $c^2 = a^2 + b^2$. Comme
$b > 0$, $c^2 > a^2$, donc $c > a$ (les longueurs sont positives) ; de même
$c > b$.

\textbf{2.} Pour $n > 2$, on a $n - 2 \geqslant 1$ et $0 < a < c$, donc
$a^{n-2} < c^{\,n-2}$. En multipliant par $a^2 > 0$ :
$a^n < a^2 c^{\,n-2}$. De même $b^n < b^2 c^{\,n-2}$.

\textbf{3.} En additionnant :
\[
a^n + b^n < (a^2 + b^2)\,c^{\,n-2} = c^2 \cdot c^{\,n-2} = c^n .
\]

\textbf{4.} Pour $n = 2$, c'est l'égalité de Pythagore $c^2 = a^2 + b^2$.
Pour $n = 1$, l'inégalité triangulaire donne au contraire $c < a + b$.
\hfill $\blacksquare$
